\section{Оптика}
\subsection{Плоское зеркало}

%%%%%%%%%%%%%%%%%%%%%%%%%%%%%%%%%%%%%%%%%%%%%%%%%%%%
%%%%%%%%%%%%%%%%%%%%%%%%%%%%%%%%%%%%%%%%%%%%%%%%%%%%
 \z{Какой минимальной высоты нужно взять зеркало, чтобы человек высотой $H$ смог увидеть себя во весь рост?}
%
{$\frac{H}{2}$}

%%%%%%%%%%%%%%%%%%%%%%%%%%%%%%%%%%%%%%%%%%%%%%%%%%%%
 \z{Под каким углом свет падает на плоскую поверхность стекла, если отражённый и преломлённый лучи образуют между собой прямой угол?}
%
{$\tan\alpha = n$}

%%%%%%%%%%%%%%%%%%%%%%%%%%%%%%%%%%%%%%%%%%%%%%%%%%%%
\z{Солнечный луч составляет с поверхностью Земли угол $40\degree$. Под каким углом к горизонту следует расположить плоское зеркало, чтобы солнечный луч попадал на дно глубокого колодца?
}
%
{$25\degree$ или $65\degree$}

%%%%%%%%%%%%%%%%%%%%%%%%%%%%%%%%%%%%%%%%%%%%%%%%%%%%
 \z[{\begin{wrapfigure}[5]{r}{0.15\textwidth}%
		\vspace{-15pt}%
		\begin{flushright}%
 			\includegraphics[width=0.15\textwidth]{4.1.4}%
		\end{flushright}%
	\end{wrapfigure}%
	}]%
{Точка $S$ движется со скоростью $v = 3\,\text{см}/\text{с}$, а зеркало — со скоростью $v_1 = 2\,\text{см}/\text{с}$, движение поступательное. С какой скоростью движется изображение точки $S$?}
%
{$u = 5\,\text{м}/\text{с}$}

%%%%%%%%%%%%%%%%%%%%%%%%%%%%%%%%%%%%%%%%%%%%%%%%%%%%
 \z{На какую высоту должен подняться в ясный солнечный день воздушный шар диаметром $10\,\text{м}$, чтобы он не отбрасывал тени на землю? Считать, что Солнце находится в зените.}
%
{$H = R_{\text{орбиты}} \,\frac{D_{\text{шара}}}{D_{\text{Солнца}}} = 1070\,\text{м}$}

%%%%%%%%%%%%%%%%%%%%%%%%%%%%%%%%%%%%%%%%%%%%%%%%%%%%
 \z[{\begin{wrapfigure}[5]{r}{0.25\textwidth}%
		\vspace{-15pt}%
		\begin{flushright}%
 			
\includegraphics[width=0.25\textwidth]{4.1.6}%
		\end{flushright}%
	\end{wrapfigure}%
	}]%
{Определите с помощью построения, куда нужно направить луч из точки $A$, чтобы после двух отражений он попал в точку $B$.}
%
{фыр}

%%%%%%%%%%%%%%%%%%%%%%%%%%%%%%%%%%%%%%%%%%%%%%%%%%%%
 \z{Точечный источник света находится между двумя плоскими зеркалами, образующими двугранный угол $\varphi$, на одинаковом расстоянии от каждого из них. Чему равен угол $\varphi$, если:
 	\pu число изображений $n = 7$;
  	\pu число изображений $n = 8$;
  	 \pu число изображений $n$?
}
%
{\puan $45\degree;$ \quad \puan  $40\degree;$ \quad \puan $\frac{360\degree}{n+1}$}



